\documentclass[11pt]{article}

\usepackage[utf8]{inputenc}
\usepackage[T1]{fontenc}
\usepackage{lmodern}
\usepackage{microtype}
\usepackage{geometry}
\usepackage{setspace}
\usepackage{hyperref}
\usepackage{graphicx}
\usepackage{booktabs}
\usepackage{enumitem}
\usepackage{amsmath}
\usepackage{amssymb}
\usepackage{float}

\geometry{margin=1in}
\setstretch{1.08}

\title{When High Arithmetic Accuracy Is Not Enough:\\
A Diagnostic Framework for Structural Robustness in Neural Arithmetic}
\author{Mohamed Mhamdi \\ \textit{Independent Researcher}}
\date{}

\begin{document}

\maketitle

\begin{abstract}
High arithmetic accuracy in neural models does not by itself establish robust reasoning. A model may perform strongly on standard evaluations while remaining fragile under structured stress. We introduce a diagnostic framework for neural arithmetic that distinguishes among distribution-bound fit, bounded compositional competence, and stronger algorithm-like behavior through a multi-dimensional scorecard. The framework combines in-distribution testing, structured adversarial families, bounded length extrapolation analysis, mechanism-sensitive probes, repeated-run validation, and qualification-aware interpretation. Using this framework, we obtain a stable baseline matrix across MLP, LSTM, and Transformer models and identify a clear architecture-dependent split on a block-boundary stress pattern. We then evaluate a first intervention based on adversarial training and show that it can produce strong gains on a seen adversarial family without yielding broad held-out robustness transfer. These results show that intervention success can be real yet still structurally narrow. More broadly, the paper argues that arithmetic evaluation should move beyond score-only reporting and toward transfer-sensitive, structure-aware diagnostics.
\end{abstract}

\section{Introduction}

High arithmetic accuracy is not, by itself, evidence of robust reasoning. Neural models can achieve strikingly strong performance on arithmetic tasks while still failing under structured stress, making standard benchmark scores insufficient for stronger claims about what the model has actually learned.

Arithmetic is an especially useful domain for exposing this problem because it has a known underlying rule structure. That makes it possible to separate two situations that are often conflated in model evaluation: a model may learn a solution that works well over a substantial but limited region of the data distribution, or it may acquire a more robust internal structure that survives distributional, structural, and procedural variation. Under standard testing conditions, these two forms of success can look deceptively similar.

This work grows out of that tension. Earlier stages of the broader research line found strong arithmetic performance under standard evaluation, but later structured tests and a full trust-recovery audit showed that performance, interpretation, and scientific confidence had to be separated much more carefully. That process restored substantial trust in the empirical backbone of the project, but it also clarified a broader lesson: high arithmetic accuracy should not be interpreted in isolation. It must be tested against structured stress and interpreted through a framework capable of distinguishing local success from broader robustness.

Project 4 was designed as the methodological answer to this problem. Instead of continuing to chase stronger raw scores, it asks a more diagnostic question: how can we distinguish between distribution-bound fit, bounded compositional competence, and stronger algorithm-like behavior in neural arithmetic models? To answer that, we introduce a diagnostic framework organized around a multi-dimensional scorecard. The framework combines in-distribution accuracy, structured adversarial evaluation, bounded length behavior, and mechanism-sensitive probes such as rounding dependence and carry corruption sensitivity. Crucially, it also includes an explicit validation layer so that interesting single-run results are not mistaken for stable findings.

Using this framework, we obtain a stable baseline matrix across three model families—MLP, LSTM, and Transformer—and identify a clear architecture-dependent split on a block-boundary stress pattern. We then test a first intervention based on adversarial training and show that it can produce a strong gain on a seen adversarial family without yielding broad held-out robustness transfer. This illustrates one of the central points of the framework: intervention success can be real and stable while still remaining structurally narrow.

The contribution of this work is therefore not the claim that a new model has solved arithmetic reasoning. Rather, the contribution is methodological. We provide a practical and reproducible way to evaluate whether apparent arithmetic progress reflects broader structural robustness or only localized gains. In doing so, we argue for a shift in arithmetic-model evaluation: away from score-only reporting and toward transfer-sensitive, structure-aware diagnostic testing.

\paragraph{Contributions.}
Our main contributions are:
\begin{enumerate}[leftmargin=1.5em]
    \item We introduce a diagnostic framework for neural arithmetic reasoning that distinguishes distribution-bound fit, bounded compositional competence, and stronger algorithm-like behavior.
    \item We define a practical multi-dimensional evaluation scorecard that combines in-distribution performance, structured adversarial testing, bounded length behavior, and mechanism-sensitive probes.
    \item We produce a stable baseline matrix across multiple model families and show that architecture-dependent differences can emerge sharply on structured adversarial patterns.
    \item We show that adversarial training can improve performance on a seen adversarial family without producing broad held-out robustness transfer.
    \item We argue that high arithmetic accuracy alone is not sufficient evidence of robust reasoning, and that arithmetic evaluation should move toward structured diagnostic testing rather than benchmark score interpretation alone.
\end{enumerate}

\section{Related Work}

This work sits at the intersection of four research threads:
\begin{enumerate}[leftmargin=1.5em]
    \item neural arithmetic and algorithm learning
    \item compositional generalization
    \item adversarial and structure-sensitive evaluation
    \item mechanistic and diagnostic analysis of model behavior
\end{enumerate}

\subsection{Neural Arithmetic and Learned Computation}

A long-standing line of work in neural networks has asked whether models can learn arithmetic procedures in a way that generalizes beyond narrow training conditions. Arithmetic has often been treated as a useful testbed for this question because it combines exact correctness conditions, explicit local and global structure, and a clear distinction between rule-following and partial approximation.

In many settings, neural models achieve strong observed performance on arithmetic tasks, especially under in-distribution or moderately shifted evaluation. However, these results leave open a deeper question: does strong performance reflect robust learned structure, or does it reflect a narrower solution that is only reliable inside a restricted region of the task space?

Project 4 contributes to this line of work by shifting emphasis away from raw arithmetic success alone and toward diagnostic characterization of what kind of success is present.

\subsection{Compositional Generalization}

The broader literature on compositional generalization has repeatedly shown that success on familiar distributions does not guarantee robust behavior under recombination, extrapolation, or structured novelty. Models may appear strong on held-out data while still relying on local statistical regularities that fail under more systematically organized variation.

This concern is especially relevant in arithmetic because arithmetic tasks naturally invite stronger interpretations. If a model performs well over many examples or longer sequences, it is easy to infer that it has learned something rule-like. But compositional-generalization research cautions against that inference unless robustness survives meaningful structural shifts.

Project 4 contributes to this area by making the distinction operational. Rather than asking only whether arithmetic behavior generalizes, it asks whether that generalization is distribution-bound, bounded but meaningful, or stronger and more structurally stable.

\subsection{Adversarial and Structure-Sensitive Evaluation}

A growing body of work in machine learning has shown that standard evaluation can fail to detect brittle or shortcut-based behavior. This has motivated the use of adversarial examples, worst-case analysis, and challenge sets in domains such as vision, language, and reasoning.

However, many adversarial approaches focus on perturbation at the input level without always preserving the internal task structure that makes the evaluation scientifically interpretable. In contrast, arithmetic allows the design of structured adversarial families whose difficulty is tied to meaningful internal task regularities rather than arbitrary perturbation.

This is one of the central motivations for Project 4. The framework uses structured adversarial arithmetic families—not just generic difficulty—to probe whether model competence transfers across qualitatively different stress conditions.

\subsection{Evaluation Beyond Aggregate Accuracy}

Recent work across machine learning has increasingly emphasized that aggregate metrics can conceal important failure modes. Average performance can hide subgroup collapse, worst-case failures, shortcut dependence, or sensitivity to small but meaningful structural changes.

This motivates richer evaluation frameworks that preserve:
\begin{itemize}[leftmargin=1.5em]
    \item per-condition reporting,
    \item worst-case behavior,
    \item and stability under reruns.
\end{itemize}

Project 4 is aligned with this direction. Its scorecard does not treat average performance as sufficient. Instead, it requires pattern-wise reporting, worst-case pattern accuracy, repeated-run validation, and multi-dimensional interpretation.

\subsection{Mechanistic and Diagnostic Analysis}

A separate but related line of work asks not only whether models succeed or fail, but what internal processes support that behavior. Mechanistic interpretability, probing, and circuit analysis all belong to this broader effort to move from performance description toward structured behavioral understanding.

Project 4 does not itself provide a full mechanistic account. Its contribution is earlier in the chain: it provides a framework that can identify the kinds of stable behavioral contrasts that would make mechanistic analysis scientifically meaningful.

\subsection{Position of This Work}

The present work differs from much of the prior literature in one key way:

\begin{quote}
it is not primarily an arithmetic-accuracy paper, nor primarily a mechanistic paper, nor merely an adversarial-evaluation paper.
\end{quote}

Instead, it is a \textbf{diagnostic framework paper}.

Its central claim is methodological:
\begin{itemize}[leftmargin=1.5em]
    \item high arithmetic accuracy should not be interpreted alone,
    \item intervention gains should not be trusted without transfer checks,
    \item and arithmetic reasoning should be evaluated through structured, validation-aware diagnostics rather than score-only reporting.
\end{itemize}

\subsection{Summary}

The relevant literature suggests three broad lessons:
\begin{enumerate}[leftmargin=1.5em]
    \item strong observed performance can be misleading,
    \item compositional generalization requires structure-sensitive evaluation,
    \item robust scientific interpretation requires more than a single scalar benchmark.
\end{enumerate}

Project 4 builds directly on those lessons by providing a practical framework for applying them to neural arithmetic.

\section{Motivation and Framing}

The motivating problem behind this work is not merely that arithmetic is difficult for neural networks. The deeper issue is that \textbf{high performance can be epistemically misleading}. A model may achieve strong results under conventional evaluation while still lacking the kind of robustness that would justify stronger reasoning claims.

Arithmetic is especially useful for studying this problem because its underlying rule structure is explicit. This makes it possible to distinguish between two situations that often look similar under standard testing: first, a model may learn a solution that works well over a substantial but limited region of the data distribution; second, a model may acquire a more robust internal structure that survives distributional, structural, and procedural variation. Standard benchmarks often fail to separate these possibilities sharply enough.

Project 4 is framed around this distinction. Its question is not simply which model achieves the best score, but what kind of competence that score actually represents.

To address that question, the framework uses three working regimes:
\begin{itemize}[leftmargin=1.5em]
    \item \textbf{Distribution-bound fit}: strong performance tied to a relatively narrow evaluation regime, with sharp degradation under more structured shifts.
    \item \textbf{Bounded compositional competence}: meaningful generalization within a bounded range, but without broader robustness to harder extrapolation or adversarial structure.
    \item \textbf{Stronger algorithm-like behavior}: more stable performance across length, adversarial structure, and equivalent formulations of the same underlying rule.
\end{itemize}

These are diagnostic categories, not metaphysical claims about cognition. Their value lies in organizing evidence in a way that prevents over-interpretation.

This leads to the central methodological principle of Project 4:

\begin{quote}
No single metric is sufficient for regime classification.
\end{quote}

A model may have strong in-distribution performance but weak worst-case adversarial behavior. Another may appear robust on one structured family while collapsing on another. An intervention may improve one family while failing to transfer beyond the patterns it explicitly saw during training. In all such cases, a single scalar metric can hide the actual scientific story.

For this reason, Project 4 uses a multi-dimensional scorecard rather than a winner-takes-all benchmark. The framework combines:
\begin{itemize}[leftmargin=1.5em]
    \item in-distribution accuracy,
    \item structured adversarial performance,
    \item length behavior,
    \item rounding sensitivity,
    \item and carry corruption sensitivity.
\end{itemize}

Structured adversarial families are especially central in this framing. Standard random test sets are useful, but they do not systematically probe the internal regularities that may matter most. By contrast, adversarial carry-pattern families place the model under sharply defined structural pressures. They make it possible to distinguish broad robustness from narrow family-specific adaptation.

This matters not only for baseline evaluation, but also for interventions. An intervention can appear successful if it improves one difficult family. That does not yet imply general robustness. The improvement may remain local to the seen family and may even trade off against performance elsewhere. Project 4 therefore treats \emph{transfer across adversarial families} as a more meaningful signal than improvement on any single family in isolation.

Project 4 is built around this idea: the goal of evaluation is not merely to rank models, but to reveal what their apparent success actually consists of. Arithmetic is the domain used here, but the methodological lesson is broader: whenever a task has a known internal rule structure, strong benchmark performance should be stress-tested against structure-sensitive diagnostics before stronger interpretive claims are made.

\section{Diagnostic Framework}

Project 4 introduces a diagnostic framework for evaluating arithmetic reasoning beyond benchmark-style accuracy reporting. The framework is designed to characterize the \emph{type} of competence a model exhibits, not merely the level of its raw performance.

It has four core components:
\begin{enumerate}[leftmargin=1.5em]
    \item working regime definitions
    \item diagnostic test families
    \item a multi-dimensional scorecard
    \item a validation layer
\end{enumerate}

Together, these components define the evidential standard used throughout the project.

\subsection{Working Regimes}

The framework organizes model behavior into three working regimes.

\paragraph{Regime 1 — Distribution-Bound Fit}
A model performs well within relatively narrow evaluation conditions, but degrades sharply under structured or mechanism-relevant shifts.

\paragraph{Regime 2 — Bounded Compositional Competence}
A model exhibits nontrivial compositional behavior within a bounded range, but does not remain stable under harder extrapolation or systematically novel structure.

\paragraph{Regime 3 — Stronger Algorithm-Like Behavior}
A model remains stable across longer lengths, structured adversarial families, and equivalent formulations of the same rule.

These regime labels are intended as diagnostic categories, not as final philosophical claims about model cognition.

\subsection{Diagnostic Test Families}

To distinguish among these regimes, the framework uses multiple evaluation families rather than a single benchmark.

\paragraph{In-Distribution Random}
This family measures standard performance on random examples drawn from the same broad family as training data. It serves as a baseline sanity condition, but it is never treated as sufficient evidence of robust reasoning.

\paragraph{Structured Adversarial}
This is the most important family in Project 4.

The v1.0 adversarial core contains:
\begin{itemize}[leftmargin=1.5em]
    \item \textbf{alternating carry}
    \item \textbf{full propagation chain}
    \item \textbf{block-boundary stress}
\end{itemize}

These patterns were selected because they place sharply different demands on carry behavior and structural generalization.

\paragraph{Length Extrapolation}
This family measures how performance changes as sequences extend beyond the core training range.

\paragraph{Mechanism-Sensitive Probes}
This family includes probes such as:
\begin{itemize}[leftmargin=1.5em]
    \item rounding sensitivity
    \item carry corruption sensitivity
\end{itemize}

These are not full mechanistic proofs. Rather, they test whether performance depends strongly on specific procedural assumptions or signal channels.

\subsection{Multi-Dimensional Scorecard}

The framework records model behavior through a scorecard rather than a single scalar.

The core scorecard dimensions in Project 4 v1.0 are:
\begin{enumerate}[leftmargin=1.5em]
    \item \textbf{in-distribution accuracy}
    \item \textbf{structured adversarial accuracy}
    \item \textbf{worst-case pattern accuracy}
    \item \textbf{length extrapolation trend}
    \item \textbf{rounding sensitivity}
    \item \textbf{carry corruption sensitivity}
\end{enumerate}

The reason for this structure is simple:

\begin{quote}
No single metric is sufficient for regime classification.
\end{quote}

Project 4 therefore requires pattern-wise reporting, worst-case reporting, and multi-dimensional interpretation.

\subsection{Pattern-Wise and Worst-Case Reporting}

Project 4 explicitly avoids reducing structured adversarial behavior to a single mean value.

For every structured adversarial evaluation, the following must be reported:
\begin{itemize}[leftmargin=1.5em]
    \item pattern-wise breakdown
    \item mean adversarial accuracy
    \item worst-case pattern accuracy
\end{itemize}

This is essential because one catastrophic failure can be hidden inside an otherwise moderate or strong average.

\subsection{Length Trend as a Diagnostic Object}

Length behavior is not treated only descriptively. It must also be expressed in a computable form.

At minimum, Project 4 requires:
\begin{itemize}[leftmargin=1.5em]
    \item accuracy by tested length
    \item relative drop beyond the training range
    \item a qualitative trend label such as:
    \begin{itemize}[leftmargin=1.5em]
        \item stable
        \item gradual decline
        \item sharp collapse
    \end{itemize}
\end{itemize}

\subsection{Intervention Evaluation Logic}

The framework is also designed to evaluate interventions, not only baselines.

An intervention is not treated as successful merely because one metric improves. Instead, the framework asks:
\begin{itemize}[leftmargin=1.5em]
    \item Did the intervention improve only a seen family?
    \item Did it transfer to unseen structured conditions?
    \item Did it improve one dimension while damaging another?
    \item Did it create local gain without broader robustness?
\end{itemize}

\subsection{Validation Layer}

Project 4 includes an explicit validation layer.

Interesting results are not enough; stable results are required.

The framework therefore distinguishes between:
\begin{itemize}[leftmargin=1.5em]
    \item single-run artifacts
    \item repeat checks
    \item stability-supported findings
\end{itemize}

At MVP stage, repeated-run validation is required before a result is elevated into a real Project 4 finding.

\subsection{Rule-Guided Classification}

Although the framework computes scorecards and generates regime guidance, final regime classification is \textbf{not fully automated}.

This is intentional. A fully automatic classifier would risk:
\begin{enumerate}[leftmargin=1.5em]
    \item overfitting interpretation to threshold artifacts
    \item concealing ambiguity in borderline cases
\end{enumerate}

Project 4 therefore uses:
\begin{itemize}[leftmargin=1.5em]
    \item rule-guided indicators
    \item scorecard-based evidence
    \item and explicit human-reviewed rationale
\end{itemize}

\subsection{What the Framework Is Designed to Detect}

The framework is especially well suited to detecting situations such as:
\begin{itemize}[leftmargin=1.5em]
    \item high standard performance with low adversarial robustness
    \item strong seen-family gains without held-out transfer
    \item architecture-specific pattern differences
    \item local success that does not scale into broader structural competence
\end{itemize}

\subsection{Summary of Framework Role}

The Project 4 framework should therefore be understood as:
\begin{itemize}[leftmargin=1.5em]
    \item a structured evidence system
    \item a robustness diagnosis tool
    \item and an anti-overclaim mechanism
\end{itemize}

Its purpose is not to decide in advance that a model is good or bad. Its purpose is to force the project to answer a harder and more useful question:

\begin{quote}
What kind of arithmetic competence is actually present, and how far does it transfer under structured stress?
\end{quote}

\section{Experimental Setup}

This section describes the empirical setup used in Project 4. The setup was designed to support the diagnostic framework rather than to maximize benchmark performance in isolation.

The experimental design has three major components:
\begin{enumerate}[leftmargin=1.5em]
    \item baseline model families
    \item diagnostic evaluation families
    \item intervention protocol
\end{enumerate}

All reported Project 4 results were interpreted under the framework's validation and qualification rules.

\subsection{Baseline Model Families}

Project 4 began with three baseline model families:
\begin{itemize}[leftmargin=1.5em]
    \item \textbf{MLP}
    \item \textbf{LSTM}
    \item \textbf{Transformer}
\end{itemize}

These baselines were selected because earlier stages of the research line had already established them as important reference families for arithmetic behavior, and because they provide meaningful variation in inductive bias.

The role of the baseline matrix was not simply to rank models. It was to establish a stable comparative reference under a common diagnostic framework.

At the MVP stage, these baselines were evaluated using a bounded but repeated-run protocol so that cross-model observations would not depend on a single accidental run.

\subsection{Baseline Source and Runtime Context}

The baseline execution paths in Project 4 were built on top of the previously audited codebase. In particular, Project 4 baseline work used bounded evaluation paths derived from the post-audit validated code context, especially the multidigit arithmetic infrastructure associated with the Phase 30 family.

These should therefore be interpreted as:
\begin{itemize}[leftmargin=1.5em]
    \item post-audit diagnostic baselines
    \item rather than raw unaudited historical reruns
\end{itemize}

\subsection{In-Distribution Evaluation}

The in-distribution family served as a baseline sanity condition. Its purpose was narrow:
- to determine whether a model performs nontrivially under a standard random evaluation path

\subsection{Structured Adversarial Evaluation}

The v1.0 adversarial family consisted of:
\begin{itemize}[leftmargin=1.5em]
    \item \textbf{alternating carry}
    \item \textbf{full propagation chain}
    \item \textbf{block-boundary stress}
\end{itemize}

These patterns probe:
- periodic carry structure
- long carry propagation dependence
- and robustness at local/global boundary interfaces

\subsection{Length Evaluation}

Project 4 also tracked bounded length-conditioned behavior in order to preserve a length-sensitive scorecard dimension without claiming broad scaling conclusions from a single curve.

\subsection{Mechanism-Sensitive Dimensions}

Two additional dimensions were included:
\begin{itemize}[leftmargin=1.5em]
    \item \textbf{rounding sensitivity}
    \item \textbf{carry corruption sensitivity}
\end{itemize}

These were useful as framework dimensions, though not equally mature across all execution paths at MVP stage.

\subsection{Validation Protocol}

Project 4 adopted a validation-aware reporting policy from the start. Baseline and intervention findings were not treated as established merely because they appeared once.

At MVP stage, repeated-run validation was applied to both:
\begin{itemize}[leftmargin=1.5em]
    \item the accepted baseline matrix
    \item and the first intervention signal
\end{itemize}

\subsection{Intervention Design}

The first Project 4 intervention was:
\begin{itemize}[leftmargin=1.5em]
    \item \textbf{adversarial training}
\end{itemize}

Its design separated:
\begin{itemize}[leftmargin=1.5em]
    \item \textbf{seen adversarial families}
    \item \textbf{held-out adversarial families}
\end{itemize}

This was necessary to distinguish:
- local intervention gain
from
- broader structural transfer

\subsection{Blockwise Extension Attempt}

Project 4 also included a first attempt at a blockwise decomposition intervention.

However, this branch was not accepted into the scientific result core because the implementation path became methodologically unresolved during debugging.

\subsection{Reporting Policy}

All Project 4 results were reported under the following rules:
\begin{enumerate}[leftmargin=1.5em]
    \item no single metric determines the interpretation
    \item pattern-wise reporting is required
    \item worst-case behavior must be preserved
    \item repeated-run stability matters
    \item regime classification remains human-reviewed and rationale-based
    \item unresolved branches are explicitly marked rather than silently dropped
\end{enumerate}

\subsection{Summary of Experimental Design}

The experimental setup of Project 4 was designed to answer a diagnostic question rather than an optimization question:
\begin{itemize}[leftmargin=1.5em]
    \item establish stable baselines
    \item expose structured failure or selectivity
    \item test whether intervention gains transfer
    \item preserve qualifications
    \item and avoid overclaiming beyond the evidence
\end{itemize}

\section{Stable Baseline Matrix}

The first major empirical outcome of Project 4 was the construction of a stable baseline matrix across three model families:
\begin{itemize}[leftmargin=1.5em]
    \item \textbf{MLP}
    \item \textbf{LSTM}
    \item \textbf{Transformer}
\end{itemize}

These baselines were not treated as one-off exploratory runs. They were repeated and checked for stability under the Project 4 validation protocol. The goal of this stage was not simply to determine which model had the highest score, but to establish whether the framework could reveal stable, structured, architecture-dependent differences.

It could.

\subsection{Why the Baseline Matrix Matters}

A diagnostic framework is only useful if it can expose meaningful structure before intervention work begins. If all baseline models look effectively identical, then the framework may not be discriminating enough. If stable cross-model differences emerge under structured evaluation, the framework is already doing something scientifically useful.

That is what happened in Project 4.

Across repeated runs, the baseline matrix showed that the three model families did \textbf{not} merely differ by small random variation. Instead, they displayed:
\begin{itemize}[leftmargin=1.5em]
    \item broadly weak exact-match performance,
    \item common failures on some structured adversarial families,
    \item and a sharp architecture-dependent split on one specific pattern family.
\end{itemize}

\subsection{Shared Weaknesses Across Baselines}

Under the current bounded Project 4 evaluation path, all three baseline families remained weak in exact-match terms.

More importantly, all three collapsed on at least some structured adversarial families, especially:
\begin{itemize}[leftmargin=1.5em]
    \item \textbf{alternating carry}
    \item \textbf{full propagation chain}
\end{itemize}

This is significant because these are exactly the kinds of structured tests that ordinary random evaluation tends to under-emphasize. Their shared collapse suggests that strong arithmetic competence should not be inferred from standard random-style success alone.

In other words, the framework immediately revealed that:
\begin{itemize}[leftmargin=1.5em]
    \item the baseline models were not broadly robust,
    \item and that structured stress testing is necessary before intervention analysis even begins.
\end{itemize}

\subsection{Architecture-Dependent Split on Block-Boundary Stress}

The strongest and most informative baseline result was the behavior on \textbf{block-boundary stress}.

This pattern produced a clear and repeated split:
\begin{itemize}[leftmargin=1.5em]
    \item \textbf{MLP}: stable success
    \item \textbf{Transformer}: stable success
    \item \textbf{LSTM}: stable failure
\end{itemize}

This was not a one-off anomaly. The pattern remained stable under repeated-run validation.

This result is scientifically important because it shows that structured adversarial evaluation can reveal differences that are not obvious from aggregate performance alone. If evaluation stopped at a single overall score, this architecture-dependent split might remain hidden.

It also shows that the baseline families differ not only in degree of weakness, but also in \textbf{type of vulnerability}. Under the Project 4 framework, the LSTM did not merely perform slightly worse than the others. It failed specifically on a pattern family where MLP and Transformer remained stable.

That makes block-boundary stress the first especially discriminative pattern family in Project 4.

\subsection{More Than a Numerical Ranking}

The baseline matrix is valuable not because it lets us say "model A is best" in a simplistic sense. Its value is diagnostic.

The baseline findings already show:
\begin{itemize}[leftmargin=1.5em]
    \item where all models fail together
    \item where models remain weak in common
    \item and where architecture-specific structure appears
\end{itemize}

This is stronger than a ranking alone because it helps identify:
\begin{itemize}[leftmargin=1.5em]
    \item which pattern families are globally difficult
    \item which pattern families are selectively difficult
    \item and which families are likely to be especially informative for intervention design
\end{itemize}

That is precisely the kind of information a diagnostic framework should surface.

\subsection{Bounded Interpretation}

The correct interpretation of the baseline matrix remains bounded.

Project 4 does \textbf{not} claim at this stage that:
\begin{itemize}[leftmargin=1.5em]
    \item any model has reached stronger algorithm-like behavior
    \item the baseline matrix settles all architecture questions
    \item or the observed differences already prove a deep mechanistic explanation
\end{itemize}

What it does support is narrower and more defensible:

\begin{quote}
the Project 4 baseline framework can detect stable architecture-dependent differences under structured stress, and those differences are meaningful enough to guide intervention design.
\end{quote}

This is already a significant methodological success.

\subsection{Why the Baseline Matrix Matters for Intervention Design}

The stable baseline matrix gave Project 4 exactly what it needed before intervention work:
\begin{itemize}[leftmargin=1.5em]
    \item a stable reference set
    \item a map of globally difficult families
    \item and a discriminative pattern family (\textbf{block-boundary stress}) for testing transfer and architecture sensitivity
\end{itemize}

That baseline map directly motivated the first MVP intervention:
\begin{itemize}[leftmargin=1.5em]
    \item adversarial training
\end{itemize}

\section{Adversarial Training as MVP Intervention}

After establishing a stable baseline matrix, Project 4 moved to its first intervention stage:
\begin{itemize}[leftmargin=1.5em]
    \item \textbf{adversarial training}
\end{itemize}

This intervention was chosen because it directly tests one of the framework's central questions:

\begin{quote}
If a model is exposed to difficult structured adversarial families during training, does that exposure produce broader robustness, or only narrow performance gains on the trained families themselves?
\end{quote}

This is exactly the kind of distinction that benchmark-driven evaluation often obscures. A model can improve after intervention and still fail to generalize in the stronger sense that matters most.

\subsection{Why Adversarial Training Was the Right First Intervention}

Adversarial training was an ideal MVP intervention because it is simple, intuitive, and directly diagnostic.

A natural expectation might be:
\begin{itemize}[leftmargin=1.5em]
    \item if a model is explicitly trained on difficult structured cases,
    \item then it should become more robust in a broader sense.
\end{itemize}

But another possibility is equally plausible:
\begin{itemize}[leftmargin=1.5em]
    \item the model may simply learn to handle the specific adversarial family it has seen,
    \item without gaining broader structural robustness.
\end{itemize}

Project 4 was designed specifically to separate these two possibilities.

This is why the intervention setup distinguished between:
\begin{itemize}[leftmargin=1.5em]
    \item \textbf{seen adversarial families}
    \item \textbf{held-out adversarial families}
\end{itemize}

Without that distinction, intervention success could easily be overstated.

\subsection{Seen-Family Gain Versus Held-Out Transfer}

The first intervention produced a clear and stable pattern:
\begin{itemize}[leftmargin=1.5em]
    \item strong improvement on a \textbf{seen adversarial family}
    \item no broad gain across all difficult families
    \item failure on a \textbf{held-out adversarial family}
\end{itemize}

This matters because it breaks a common but weak inference:
\begin{itemize}[leftmargin=1.5em]
    \item "the intervention improved performance, therefore it improved robustness"
\end{itemize}

Project 4 shows that this inference is too crude.

The intervention produced a real effect, but the effect was \textbf{narrow} rather than broadly transferable.

\subsection{Why This Is a Strong Result}

This result is important not because it shows adversarial training "failed" in a simplistic sense. The intervention did succeed in changing the model's behavior.

What makes the result scientifically valuable is that the framework can distinguish:
\begin{itemize}[leftmargin=1.5em]
    \item \textbf{real gain}
    \item from
    \item \textbf{robust transfer}
\end{itemize}

The first intervention therefore validates the framework itself.

If Project 4 had reported only:
\begin{itemize}[leftmargin=1.5em]
    \item one improved adversarial score
\end{itemize}

then the intervention might look broadly successful.
But because the framework includes:
\begin{itemize}[leftmargin=1.5em]
    \item held-out families
    \item repeated-run stability
    \item pattern-wise reporting
    \item and explicit qualification
\end{itemize}

the interpretation becomes much sharper.

This is exactly the kind of result the framework was designed to produce.

\subsection{Main Diagnostic Conclusion}

The most defensible interpretation of the first intervention is:

\begin{quote}
adversarial training can improve performance on a specifically seen structured family without producing broad structural robustness transfer.
\end{quote}

This is not a trivial result.

It tells us that:
\begin{itemize}[leftmargin=1.5em]
    \item intervention gains should not be treated as automatically general
    \item structured evaluation families must remain distinct
    \item and held-out structured stress is necessary for meaningful robustness evaluation
\end{itemize}

In that sense, the intervention confirmed the need for the framework.

\subsection{Why This Result Is Methodologically Strong}

The intervention result is scientifically strong because it is:
\begin{itemize}[leftmargin=1.5em]
    \item repeated
    \item stable
    \item structurally interpretable
    \item and framed against a baseline matrix
\end{itemize}

That allows a more precise conclusion than:
\begin{itemize}[leftmargin=1.5em]
    \item "the intervention helped"
    \item or
    \item "the intervention failed"
\end{itemize}

Instead, the result shows:
\begin{itemize}[leftmargin=1.5em]
    \item the intervention produced a real gain
    \item the gain was family-specific
    \item and broad transfer did not follow automatically
\end{itemize}

That is a much more useful scientific conclusion.

\subsection{Bounded Interpretation}

The interpretation remains intentionally bounded.

Project 4 does \textbf{not} claim from this intervention alone that:
\begin{itemize}[leftmargin=1.5em]
    \item adversarial training can never produce broader robustness
    \item all intervention gains are merely memorization
    \item or the exact internal cause of the observed trade-off has been identified
\end{itemize}

Those claims would be too strong.

What the current evidence supports is narrower:
\begin{itemize}[leftmargin=1.5em]
    \item the first intervention yielded stable narrow gain
    \item broad held-out transfer was not observed
    \item and intervention behavior must therefore be evaluated family-by-family rather than by aggregate improvement alone
\end{itemize}

\subsection{Why This Matters Beyond This Intervention}

The significance of this result extends beyond the specific adversarial-training setup used here.

It supports a broader methodological lesson:

\begin{quote}
intervention success should be judged by transfer across structured stress families, not by local gain alone.
\end{quote}

That lesson could matter well beyond arithmetic. Wherever a model is evaluated on:
\begin{itemize}[leftmargin=1.5em]
    \item compositional structure
    \item adversarially organized inputs
    \item or distribution-sensitive tasks
\end{itemize}

there is a risk that local gains will be mistaken for broader capability.

Project 4 provides one concrete framework for avoiding that mistake.

\subsection{Why This Was Enough for MVP}

At this point in the project, the first intervention had already succeeded in one important sense:
\begin{itemize}[leftmargin=1.5em]
    \item it showed that the framework can separate local improvement from broader robustness transfer
\end{itemize}

That was enough for MVP-level scientific value.

This is why the accepted Project 4 synthesis rests primarily on:
\begin{itemize}[leftmargin=1.5em]
    \item the stable baseline matrix
    \item and the stable first intervention signal
\end{itemize}
rather than on unresolved extension branches.

\section{What the Framework Reveals}

The value of Project 4 does not lie only in the individual numbers reported for baselines or interventions. Its real value lies in what becomes visible once those numbers are placed inside a diagnostic structure.

What the framework reveals is not simply which model is better, or whether one intervention improved a metric. It reveals the \emph{shape} of arithmetic competence and the \emph{limits} of that competence under structured stress.

This is the central contribution of Project 4.

\subsection{Why Standard Accuracy Is Not Enough}

If evaluation stopped at in-distribution accuracy alone, the scientific picture would be severely incomplete.

A model could appear competitive under standard random evaluation while still:
\begin{itemize}[leftmargin=1.5em]
    \item collapsing on particular structured families
    \item depending on narrow support conditions
    \item or improving only on the exact families it saw during intervention
\end{itemize}

The framework reveals that these are often the most informative parts of the result.

Standard accuracy tells us whether a model performs.
The framework asks \emph{what kind of performance that is}.

\subsection{Why Family Identity Matters}

One of the clearest lessons from the baseline and intervention results is that arithmetic weakness is not uniform.

Different pattern families expose different forms of limitation.

A model can:
\begin{itemize}[leftmargin=1.5em]
    \item fail catastrophically on one structured family
    \item remain strong on another
    \item and improve on one family while not transferring to another
\end{itemize}

Without explicit family-level structure, these differences would collapse into an average and become harder to interpret.

The framework reveals that \textbf{family identity matters}.

That is scientifically important because it shifts the evaluation question from:
\begin{itemize}[leftmargin=1.5em]
    \item "Did accuracy go up?"
\end{itemize}
to
\begin{itemize}[leftmargin=1.5em]
    \item "On what structure did it go up, and did that transfer?"
\end{itemize}

\subsection{Why Baselines and Interventions Must Be Read Together}

Another key lesson is that baselines and interventions cannot be interpreted in isolation.

The stable baseline matrix tells us:
\begin{itemize}[leftmargin=1.5em]
    \item where models already differ before intervention
    \item which families are globally hard
    \item and which families are especially discriminative
\end{itemize}

The intervention result then tells us:
\begin{itemize}[leftmargin=1.5em]
    \item whether intervention changes those baseline patterns
    \item whether it improves one family only
    \item or whether it shifts robustness in a broader way
\end{itemize}

This combined reading is stronger than either baselines or interventions alone.

Project 4 therefore reveals not only behavior, but \textbf{change relative to a structured baseline map}.

\subsection{Stable Weakness Is Scientifically Valuable}

A stable failure pattern is often more informative than a high score with no diagnostic structure.

A model that repeatedly fails on one adversarial family but succeeds on another tells us more than a single average would. Likewise, an intervention that repeatedly improves a seen family while failing to transfer still yields a scientifically meaningful result.

This is one of the strongest commitments of Project 4:
\begin{itemize}[leftmargin=1.5em]
    \item negative and qualified results are not consolation prizes
    \item they are central evidence
\end{itemize}

\subsection{Why This Matters for Arithmetic Research}

Arithmetic is often treated as a benchmark domain because it has a known rule structure and a clear notion of correctness. But that same clarity can create false confidence if evaluation is too narrow.

Project 4 shows that arithmetic reasoning research should be careful in at least three ways:
\begin{enumerate}[leftmargin=1.5em]
    \item benchmark success should not be interpreted alone
    \item structured adversarial evaluation is not optional
    \item intervention claims require transfer testing
\end{enumerate}

These are methodological lessons, but they are also scientific ones.

\subsection{Why This Matters Beyond Arithmetic}

The broader relevance of the framework comes from the fact that arithmetic is only one domain where benchmark scores can conceal narrow solutions.

The same problem appears whenever a task has:
\begin{itemize}[leftmargin=1.5em]
    \item latent compositional structure
    \item multiple possible shortcut strategies
    \item or a difference between standard examples and structured stress cases
\end{itemize}

The broader methodological logic may therefore generalize:
\begin{itemize}[leftmargin=1.5em]
    \item evaluate family-by-family
    \item preserve worst-case structure
    \item validate repeated runs
    \item separate local gain from transfer
\end{itemize}

The exact scorecard will still need domain-specific redesign, but the underlying logic is portable.

\subsection{The Main Methodological Takeaway}

The strongest lesson that Project 4 makes explicit is:

\begin{quote}
A model can become better in a real sense without becoming broadly robust in the sense that matters most.
\end{quote}

This is why Project 4 is a diagnostic framework project rather than a benchmark race.

Its value lies in forcing a stronger standard of evidence:
\begin{itemize}[leftmargin=1.5em]
    \item not "did the score improve?"
    \item but "what exactly improved, where, and how far did that improvement transfer?"
\end{itemize}

\subsection{Why the MVP Is Already Enough}

Project 4 does not need a Regime 3 model to matter.

The framework already succeeded because it produced:
\begin{itemize}[leftmargin=1.5em]
    \item a stable baseline matrix
    \item a stable intervention signal
    \item and a principled distinction between narrow gain and broader transfer
\end{itemize}

That is enough to justify the MVP as a meaningful scientific contribution.

The framework revealed something important:
\begin{itemize}[leftmargin=1.5em]
    \item improvement can be genuine
    \item and still remain structurally limited
\end{itemize}

That is the key result Project 4 contributes to the research line.

\section{Limitations}

Although Project 4 achieved its MVP objective, its conclusions remain intentionally bounded. This is a strength rather than a weakness: the framework was designed to separate what the evidence supports from what it does not yet support.

Several limitations therefore remain explicit.

\subsection{No Regime 3 Claim}

Project 4 does not establish that any tested model reached stronger algorithm-like behavior.

This was never required for MVP success, and no such claim should be inferred from the current results.

\subsection{Uneven Maturity Across Scorecard Dimensions}

The framework includes multiple dimensions, but not all are equally mature in implementation at the current stage.

In particular:
\begin{itemize}[leftmargin=1.5em]
    \item some dimensions were already central and empirically informative, especially structured adversarial behavior and stable baseline/intervention comparisons
    \item other dimensions, such as some mechanism-sensitive probes, remain less mature in bounded implementation form
\end{itemize}

This does not invalidate the framework, but it means the framework should be read as:
\begin{itemize}[leftmargin=1.5em]
    \item operational and useful
    \item but still evolving in implementation maturity
\end{itemize}

\subsection{Blockwise Decomposition Remains Unresolved}

Project 4 included a first blockwise decomposition attempt, but that branch was not accepted into the scientific result core because the implementation path became methodologically unresolved during debugging.

This means:
\begin{itemize}[leftmargin=1.5em]
    \item Project 4 does not yet provide a clean answer to the blockwise structural hypothesis
    \item and stronger claims about chunking or carry-interface decomposition remain open
\end{itemize}

That branch is documented, but excluded from the accepted MVP result core.

\subsection{No Full Mechanistic Proof}

Project 4 is a diagnostic framework and intervention study, not a completed mechanistic interpretability analysis.

It can reveal:
\begin{itemize}[leftmargin=1.5em]
    \item stable weaknesses
    \item transfer failures
    \item architecture-dependent differences
    \item and family-specific intervention gains
\end{itemize}

But it does not yet prove:
\begin{itemize}[leftmargin=1.5em]
    \item the exact internal computation strategy of any model
    \item the representational basis of carry handling
    \item or the circuit-level reason why one architecture succeeds where another fails
\end{itemize}

Those remain valid future directions.

\subsection{Bounded Scope}

The implementations in Project 4 were designed to be operational, reproducible, and bounded rather than maximally exhaustive.

This means:
\begin{itemize}[leftmargin=1.5em]
    \item not every architecture or intervention variant was explored
    \item some paths were intentionally simplified at MVP stage
    \item engineering completeness was sometimes deferred in favor of a stable diagnostic core
\end{itemize}

This is acceptable for MVP, but it limits how far the current conclusions should be stretched.

\subsection{Arithmetic as a Special Domain}

Arithmetic is unusually useful for this kind of work because:
\begin{itemize}[leftmargin=1.5em]
    \item correctness is explicit
    \item structure is well-defined
    \item and stress patterns can be designed cleanly
\end{itemize}

But this also means caution is needed when generalizing outward.

Project 4 supports a methodological claim that likely extends beyond arithmetic:
\begin{itemize}[leftmargin=1.5em]
    \item benchmark success should be stress-tested against structured evaluation
\end{itemize}

However, it does not by itself prove that the exact same families, thresholds, or regime behavior will transfer unchanged to other domains.

\subsection{Diagnosis Rather Than Final Theory}

Project 4 is strongest when read as:
\begin{itemize}[leftmargin=1.5em]
    \item a diagnostic system
    \item a measurement philosophy
    \item and a structured intervention-testing framework
\end{itemize}

It is weaker if over-read as:
\begin{itemize}[leftmargin=1.5em]
    \item a final theory of arithmetic reasoning in neural networks
    \item a definitive architecture ranking
    \item or a complete mechanistic account
\end{itemize}

This is not a defect. It is part of what makes the project's claims credible.

\subsection{Why These Limitations Matter}

These limitations are not decorative disclaimers. They define the boundary of the project's valid contribution.

Project 4 established:
\begin{itemize}[leftmargin=1.5em]
    \item that structured diagnostics can reveal baseline differences hidden by simpler evaluation
    \item that intervention gains can be narrow and non-transferable
    \item and that repeated-run stability matters for interpreting those signals
\end{itemize}

That is already substantial.

The limitations simply prevent those findings from being overstated.

\subsection{Practical Interpretation Boundary}

The correct final use of Project 4 is therefore:
\begin{itemize}[leftmargin=1.5em]
    \item as a meaningful MVP diagnostic contribution
    \item as a stable baseline/intervention study
    \item and as a framework that can be extended later
\end{itemize}

It should not yet be used as proof that the full problem of arithmetic robustness has been solved, nor as a complete account of the mechanisms underlying the observed patterns.

\section{Conclusion}

This work began from a simple but increasingly important problem: high arithmetic accuracy is not enough to justify strong claims about robust reasoning. Neural models can appear highly successful under standard evaluation while remaining fragile under structured stress. If evaluation stops at benchmark-style performance, this distinction can remain hidden.

Project 4 addressed this problem by introducing a diagnostic framework for arithmetic reasoning. Rather than asking only whether a model performs well, the framework asks what kind of competence that performance represents. To do this, it combines a multi-dimensional scorecard, structured adversarial families, repeated-run validation, and rule-guided interpretation. The goal is not merely to rank models, but to distinguish between distribution-bound fit, bounded compositional competence, and stronger algorithm-like behavior.

Using this framework, we obtained a stable baseline matrix across three model families and identified a clear architecture-dependent split on block-boundary stress. We then tested a first intervention based on adversarial training and found that it can produce strong gains on a seen adversarial family without yielding broad held-out robustness transfer. This is the key empirical result of the project: intervention gains can be real, stable, and still structurally narrow.

The broader contribution of Project 4 is methodological. It shows that arithmetic evaluation should move beyond raw accuracy reporting and toward transfer-sensitive, structure-aware diagnostics. In this sense, the project contributes not a final theory of arithmetic reasoning, but a more rigorous way to evaluate it.

This is why Project 4 succeeds even without a Regime 3 model. Its MVP objective was not to prove that stronger algorithm-like behavior had been found, but to build a framework capable of distinguishing local gains from broader robustness. That objective was met. The project now provides:
\begin{itemize}[leftmargin=1.5em]
    \item a working diagnostic framework,
    \item a stable baseline comparison matrix,
    \item and a stable first intervention signal with a scientifically meaningful interpretation.
\end{itemize}

At the same time, the work remains explicitly bounded. It does not provide a full mechanistic account of model behavior, it does not resolve the blockwise hypothesis, and it does not claim that the current framework exhausts every relevant diagnostic dimension. These limitations are not weaknesses to hide. They are part of what makes the current contribution credible.

The strongest final conclusion is therefore straightforward:

\begin{quote}
high arithmetic performance is not sufficient evidence of robust reasoning, and intervention gains should be evaluated by their ability to transfer across structured stress families rather than by local improvement alone.
\end{quote}

Project 4 establishes a practical framework for making that distinction. That is its central contribution, and it provides a concrete methodological foundation for future work on arithmetic reasoning, robustness, and structured generalization in neural systems.

\end{document}
